% Part III opener: the scope of Chapters 17--24, applications in science and medicine.
\noindent Part~III carries the architectures of Parts~I and~II out of the idealized
benchmark problems they were built on and into multi-scale physical and biomedical
systems, where the governing geometry is a triangulated car surface rather than a
torus, the coefficients are heterogeneous tensors, and the deliverable is a force
coefficient, a lead-time skill score, a plume contour, or an image. Its organizing
question is the multi-query bottleneck: five domains --- computational fluid
dynamics, geophysical climate modeling, geological carbon sequestration, cardiac
electrophysiology, and inverse problems in medical tomography --- were selected
because each issues the same equation many times over, and it is that repetition,
not the accuracy of any single solve, that operator learning answers. For each
domain the Part establishes the underlying partial differential equation (PDE)
formulation, analyzes continuous boundary condition enforcement, details the
architectural adaptation it required --- the Geometry-Informed Neural Operator
(GINO), FourCastNet on the Adaptive Fourier Neural Operator (AFNO), the
U-Net-enhanced Fourier Neural Operator (U-FNO), the Physics-Informed Neural
Operator (PINO), and, in tomography, a Fourier-neural-operator-class
reconstruction of undersampled data, arranged either as an unrolled
data-consistency cascade or as a learned filtering stage inside filtered
backprojection (FBP) --- and
evaluates quantitative speedup and fidelity against the classical solvers it
displaces: the finite element method (FEM), the finite volume method (FVM),
implicit finite difference, and pseudo-spectral or spectral-transform dynamics.
Accounting is treated as part of the evaluation. Every ratio is read against the
baseline, the hardware, and the accuracy target that produced it, and a multiplier
whose pair of timings cannot be recovered is recorded as unestablished rather than
printed. The Part closes with an interdisciplinary taxonomy synthesizing those
performance metrics, the physical conservation guarantees each architecture does
and does not carry, and the translational roadblocks standing between a published
surrogate and clinical or industrial deployment.
